\documentclass{standalone} % Utilise la classe standalone pour compiler uniquement le TikZ
%\include{pakage_tikz}
\usepackage{tikz}          % Charge le package TikZ
\usepackage{xcolor}        % Gestion avancée des couleurs
\usepackage{amsmath}       % Gestion des formules mathématiques
\usetikzlibrary{arrows.meta} % Pour gérer les flèches avec des formes spécifiées comme "triangle"
\usetikzlibrary {decorations.pathmorphing}
\usepackage{pgfplots}
\pgfplotsset{compat=1.18}
\usepackage[active,tightpage]{preview}
\PreviewEnvironment{tikzpicture}
\usetikzlibrary{arrows.meta}
\usetikzlibrary {arrows.meta,bending,positioning}
\usetikzlibrary {datavisualization.formats.functions}
%PREAMBULE pour schÃéma
\usepackage{pgfplots}
\usepackage{tikz}
\usepackage[european resistor, european voltage, european current]{circuitikz}
\usetikzlibrary{arrows,shapes,positioning}
\usetikzlibrary{decorations.markings,decorations.pathmorphing,
decorations.pathreplacing}
\usetikzlibrary{calc,patterns,shapes.geometric}
\usepackage{anyfontsize}

\setlength\PreviewBorder{5pt}
%\usepgfplotslibrary{external}
%\tikzexternalize

%% Automatiser cela : une solution plus élégante  nouvelle page, tu peux utiliser la commande \newpage juste avant chaque \begin{tikzpicture}
%\let\oldtikzpicture\tikzpicture
%\let\endoldtikzpicture\endtikzpicture
%\renewenvironment{tikzpicture}{\newpage\oldtikzpicture}{\endoldtikzpicture}
%\usepackage{etoolbox} % dans le préambule, si pas encore utilisé
%\pretocmd{\tikzpicture}{\clearpage}{}{}


% Définition de la fonction pour générer la grille
\newcommand{\drawgrid}[4]{


	\begin{scope}[opacity = 0.2]
	
	
    	% Paramètres : #1=xmin, #2=xmax, #3=ymin, #4=ymax

   	 	% Grille principale en cm
    	\draw[very thin, gray] (#1,#3) grid (#2,#4); % Grille de base (1 cm)
    
    	\draw[very thin, gray] (#1,#3) grid[step=0.1] (#2,#4);

    	\draw[thin, black] (#1,#3) edge[thick,line width=0.8ex,->,>=Stealth ] (#2,#3); 
    
    	\foreach \x in {#1,..., #2} {
        	\draw[shift={(\x, #3)}] (0,-0.1) edge[line width=0.8ex] node[pos = -1 , below] {\x} (0,0.1); % Lignes verticales en mm
    	}

    	\draw[thin, black] (#1,#3)edge[thick,line width=0.8ex,->,>=Stealth ] (#1,#4); 
    
    	\foreach \y in {#3,..., #4} {
        	\draw[shift={(#1, \y)}] (-0.1 , 0 ) edge[line width=0.8ex] node[pos = -1 , left ] {\y} (0.1 , 0); % Lignes verticales en mm
    	}
    
    \end{scope}


}


\newcommand{\Axex}[1][$x$]{ % "x" est la valeur par défaut de #1
    \draw
        (-10, 0) edge[thick, line width=0.8ex, ->, >=Triangle, color=\colorThree] 
        node[shift = {(0,1)} , pos=1.01, below]{ \huge #1 } (10, 0)
        ;
}

\newcommand{\Axetheta}[1][$\theta$]{
	\draw
		(-10, 0 ) edge [thick,line width=0.8ex,->,>=Triangle , color = \colorThree ]node [pos=1.01,right]{\huge #1 } ( 10  , 0 )
	;
}

\newcommand{\Axedensity}[1][$n$]{
	\draw
		(0, -1 ) edge [thick,line width=0.8ex,->,>=Triangle , color = \colorThree ]node [pos=1.01,left  ]{\huge #1 } ( 0  , 6 )
	;
}

\newcommand{\Axedistribution}[1][$\nu$]{
	\draw
		(0, -1 ) edge [thick,line width=0.8ex,->,>=Triangle , color = \colorThree ]node [pos=1.01,left  ]{\huge #1 } ( 0  , 6 )
	;
}

\newcommand{\Axesphase}[2][$x$]{
	\Axex[#1]
	\begin{scope}[rotate = 90]
		\Axetheta[#2]	
	\end{scope}
}

\newcommand{\Axesdensity}[2][$x$]{
	\Axex[#1]
	\Axedensity[#2]
}

\newcommand{\Axesdistribution}[2][\theta]{
	\Axetheta[#1]
	\Axedistribution[#2]
}

\newcommand{\Axetemps}[1][Deformation]{
	\draw 
		(-5 , 0 ) edge [line width=2.8ex, color=\colorOne ]( 0, 0)
		(43 , 0 ) edge [line width=2.8ex, color=\colorOne , ->, >=Triangle] node[pos = 1.0 , right , scale = 2 ] { \fontsize{60pt}{720pt}\selectfont $\tilde{t}$} ( 50 , 0)
		( 0 , 1 ) edge [line width=2.0ex, color=\colorOne ] ( 0 , -1 )  
		( 43 , 1 ) edge [line width=2.0ex, color=\colorOne ] ( 43 , -1 ) 
		(0 , 0) edge[<->, > = stealth, line width=2.8ex, color=\colorOne ]  node [ rectangle , midway, fill = white , scale = 2 ] {\fontsize{600pt}{720pt}\selectfont \bf  #1}(43, 0)  
	;

}

\newcommand{\Palette}[1][\colorOne]{
	\clip[decorate, decoration={random steps, segment length=3pt, amplitude=3pt}]
        	(-1,-1) rectangle (1,1);
	\draw[fill=#1 , color = #1 ] (-2,-2) rectangle (2,2);
}






% Définition des couleurs avec les codes HTML
\definecolor{colorOne}{HTML}{443E46}
%\definecolor{colorTwo}{HTML}{F6DEB8}
\definecolor{colorTwo}{HTML}{FFFFFF}

\definecolor{colorThree}{HTML}{908CA4}
\definecolor{colorFour}{HTML}{57659E}
\definecolor{colorFive}{HTML}{C57284}
\definecolor{colorSix}{HTML}{FF5B69}

% Raccourcis pour les couleurs
\def\colorOne{colorOne}
\def\colorTwo{colorTwo}
\def\colorThree{colorThree}
\def\colorFour{colorFour}
\def\colorFive{colorFive}
\def\colorSix{colorSix}



% Graduation ordonné
	
\newcommand{\GraduationYMm}[2]{
	\foreach \r in {1,...,9} {
		\draw 
			({-0.25 * 0.5 * (cos(3.14*0.2*\r r)^2 + 1)}, #1*#2 - #2 + 0.1*#2*\r) 
			edge [thick, line width=0.3ex, color=\colorOne] 
			({0.25 * 0.5 * (cos(3.14*0.2*\r r)^2 + 1)}, #1*#2 - #2 + 0.1*#2*\r);
	}
}

\newcommand{\GraduationYCm}[4]{% #1: min, #2: max, #3: liste
  \foreach \R in {#1,...,#2} {
    \pgfmathparse{#3[\R-(#1)]} % index dans la liste
    \let\labelText\pgfmathresult

    \draw[color=\colorThree]
      (-0.3, \R*#4)
      edge [thick, line width=0.5ex, color=\colorOne]
      node [pos=-0.0, left]{\Large \labelText}
      (0.3, \R*#4);

    %\GraduationYMm{\R}{#4}
  }

  % Graduation supplémentaire à #2 + 3
  \pgfmathparse{#2 + 1}
  \edef\extendedPosition{\pgfmathresult}
  %\GraduationYMm{\extendedPosition}{#4}
}

\newcommand{\GraduationXMm}[2]{
	\foreach \r in {1,...,9} {
		\draw 
			( #1*#2 - #2 + 0.1*#2*\r,{-0.25 * 0.5 * (cos(3.14*0.2*\r r)^2 + 1)}) 
			edge [thick, line width=0.3ex, color=\colorOne] 
			( #1*#2 - #2 + 0.1*#2*\r,{0.25 * 0.5 * (cos(3.14*0.2*\r r)^2 + 1)});
	}
}

\newcommand{\GraduationXCm}[4]{% #1: min, #2: max, #3: liste
  \foreach \R in {#1,...,#2} {
    \pgfmathparse{#3[\R-(#1)]} % index dans la liste
    \let\labelText\pgfmathresult

    \draw[color=\colorThree]
      (\R*#4,-0.3)
      edge [thick, line width=0.5ex, color=\colorOne]
      node [pos=-0.0, below]{\Large \labelText}
      (\R*#4,0.3);

    %\GraduationXMm{\R}{#4}
  }

  % Graduation supplémentaire à #2 + 3
  \pgfmathparse{#2 + 1}
  \edef\extendedPosition{\pgfmathresult}
  %\GraduationXMm{\extendedPosition}{#4}
}

% Définition de la fonction pour générer la grille
\newcommand{\backgrowngrid}[6]{


	\begin{scope}[xscale = #5 , yscale = #6 ]
	
	
    	% Paramètres : #1=xmin, #3=xmax, #2=ymin, #4=ymax
    	%\clip[scale = 0.95 , decorate, decoration={random steps, segment length=0.1 cm, amplitude=0.2cm}](#1 , #2 ) rectangle ( #3 , #4 ) ;
    	\pgfmathparse{#1-1}  
    	\let\valXmin\pgfmathresult
    	\pgfmathparse{#2-1}  
    	\let\valYmin\pgfmathresult
    	\pgfmathparse{#3+1}  
    	\let\valXmax\pgfmathresult
    	\pgfmathparse{#4+1}  
    	\let\valYmax\pgfmathresult
    	
    	\draw[fill = \colorThree , opacity = 0.1]  (\valXmin , \valYmin ) rectangle ( \valXmax , \valYmax ) ;

   	 	% Grille principale en cm
    	\draw[very thin, line width=0.5ex ,  color=\colorTwo] (\valXmin , \valYmin ) grid ( \valXmax , \valYmax ); % Grille de base (#5 cm , #6 cm )
    
    	%\draw[very thin, line width=0.3ex ,  color=\colorTwo] (#1,#2) grid[step=0.1] (#3,#4);

    	%\draw[thin, black] (#1,#2) edge[thick,line width=0.8ex,->,>=Stealth ] (#3,#2); 
    
    	%\foreach \x in {#1,..., #3} {
    	%	\pgfmathparse{int(\x * #5)} % calcule et arrondit \pgfmathparse{round((\x * #5)/10)*10} Par exemple, au multiple de 10 le plus proche 
    	%	\let\val\pgfmathresult
        %	\draw[shift={(\x, #2)}] (0,-0.1) edge[line width=0.8ex] node[pos = -1 , below] {\val} (0,0.1); % Lignes verticales en mm
    	%}

    	%\draw[thin, black] (#1,#2)edge[thick,line width=0.8ex,->,>=Stealth ] (#1,#4); 
    
    	%\foreach \y in {#2,..., #4} {
    	%	\pgfmathparse{int(\y * #6)} % calcule et arrondit  
    	%	\let\val\pgfmathresult
        %	\draw[shift={(#1, \y)}] (-0.1 , 0 ) edge[line width=0.8ex] node[pos = -1 , left ] {\val} (0.1 , 0); % Lignes verticales en mm
    	%}
    
    \end{scope}


}



\begin{document}

\begin{tikzpicture}
			%\input{figures/04_GGE_Fluctuation/Occupation_code}	
			\begin{scope}[transform canvas={scale=0.6}]
			\input{Occupation_theta_code}	
			\end{scope}
			
			\draw[color = red , scale = 0.5 , draw = none ] (-13.5 , -1) rectangle (13 , 10) ; 
				
			
		\end{tikzpicture}
		
\begin{tikzpicture}
\begin{axis}[
    xlabel={$\theta~(\mu\text{m}/\text{ms})$},
    ylabel={$\nu(\theta)$},
    grid=both,
    width=10cm,
    height=6cm
]
\addplot+[smooth, \colorFour , thick , no marks , line width=0.5ex] table {../02_GGE_TBA/nu_vs_theta.dat};
\end{axis}

%\draw[line width=1.0ex, color=\colorFour, fill=none]
%    plot[smooth] file {../02_GGE_TBA/nu_vs_theta.dat}
%    (6,0.8) edge[] node[shift={(0,1)}, pos=1, right]
%        {\fontsize{16pt}{19pt}\selectfont \bf $\nu(\theta)$} (6.5,0.8);
\end{tikzpicture}

\begin{tikzpicture}
\begin{axis}[
    xlabel={$\theta~(\mu\text{m}/\text{ms})$},
    ylabel={$(\text{ms}/\mu\text{m}^2)$},
    grid=both,
    width=10cm,
    height=6cm,
    legend style={at={(0.98,0.98)}, anchor=north east, font=\small}
]

x
% rho_s
\addplot+[smooth, thick, \colorSix, no marks , line width=0.5ex] table {../02_GGE_TBA/rho_s_vs_theta.dat};
\addlegendentry{$\rho_s(\theta)$}

% rho
\addplot+[smooth,\colorFour, no marks , line width=0.5ex] table {../02_GGE_TBA/rho_vs_theta.dat};
\addlegendentry{$\rho(\theta)$}

\end{axis}
\end{tikzpicture}


\begin{tikzpicture}
    
        
    % Rectangle avec texte "Tonks Girardeau"
    
        
    \fill[fill=\colorFive , ] (6,0) -- (6,2) -- (9,2) -- (9,0) ;
    \node[draw = none , fill=none, minimum width=3cm, minimum height=2cm, align=center, \colorTwo] 
        at (7.5,0.5) {\fontsize{10pt}{16pt}\selectfont Tonks Girardeau};
        
    \fill[fill=\colorSix] (6,2) -- (0,5) -- (0,6) -- (9,6) -- (9,2) ;
    \node[draw = none , fill=none, minimum width=3cm, minimum height=2cm, align=center , \colorTwo] 
        at (5,5) {\fontsize{10pt}{16pt}\selectfont Gaz de Bose idéal};
        
    \fill[fill=\colorFour] (6,2) -- (0,5) -- (0,0) -- (6,0)  ;
    \node[draw = none , fill=none, minimum width=3cm, minimum height=2cm, align=center, \colorTwo] 
        at (3,0.5) {\fontsize{10pt}{16pt}\selectfont Quasi-condensat de Bose-Einstein};
        
   \draw[line width=0.4ex, color=\colorOne , dashed ]
   		(6,0) -- (6,2) -- (9,2) 
   ; 
          
     
     \draw[line width=0.8ex, color=\colorOne] 
    	(0,5) edge[line width=0.5ex, color=\colorOne]  (6,2) 
    	(0,4) edge[line width=0.5ex, color=\colorFive] node[pos = 0.2, sloped, above] {\fontsize{6pt}{16pt}\selectfont régime thermique} node[pos = 0.2, sloped, below] {\fontsize{6pt}{16pt}\selectfont régime quantique} (6,2)
    	(0,6) edge[line width=0.5ex, color=\colorFour] node[pos = 0.2, sloped, above] {\fontsize{6pt}{16pt}\selectfont régime classique} node[pos = 0.2, sloped, below] {\fontsize{6pt}{16pt}\selectfont régime dégénéré} (9,0)
    	
    ;
    
    
    \draw[thick, color=\colorThree, fill=\colorThree , opacity = 0.8] (2.6,3.) ellipse [x radius=0.8, y radius=0.3];
    
    
    % Axe horizontal (gamma)
    \draw[ line width=0.8ex, color=\colorOne] 
        (0,0) -- (9,0) 
        node[shift = {(0,-1)} , pos=0.5, below] {\fontsize{16pt}{19pt}\selectfont $\gamma$};

    % Axe vertical (t)
    \draw[ line width=0.8ex, color=\colorOne] 
        (0,0) -- (0,6) 
        node[ shift = {(-1,0)} , pos=0.5, above, rotate=90] {\fontsize{16pt}{19pt}\selectfont $t$};
        
    % --- Définition des listes ---
	\def\listetuple{0/10^{-4}, 1.5/10^{-3} , 3/10^{-2} , 4.5/10^{-1} , 6/1 , 7.5/10 , 9/10^{2}  }
	\def\listetrais{0,1.5,3,4.5,6,7.5,9}

	% --- Boucle sur toutes les positions ---
	\foreach \r in \listetrais {
    	% Initialisation
    	\global\def\found{0}
    	\xdef\nomtheta{}
    
   	 	% Vérification si \r est dans listetuple
    	\foreach \x/\y in \listetuple {
        	\ifdim \r pt=\x pt
            	\global\def\found{1}
            	\xdef\nomtheta{\y}
        	\fi
    	}

    	% Tracé de la graduation
    	\ifnum\found=1
        	% Graduation avec étiquette
        	\draw[color=\colorOne, thick, line width=0.5ex] 
            	(\r,0.1) -- (\r,-0.1) node[shift = {(0,-0)} , \colorOne , pos=1 , below] {\large $\nomtheta$};
        	%\filldraw[line width=0.5ex, color=\colorSix] (\r,0) circle (4pt);
    		\else
        	% Graduation sans étiquette
        	\draw[color=\colorOne, thick, line width=0.5ex] 
            	(\r,-0.1) -- (\r,0.1);
        	%\filldraw[line width=0.5ex, color=\colorTwo] (\r,0) circle (4pt); 
    	\fi
	}
	
	
	% --- Définition des listes ---
	\def\listetuple{0/10^{-4}, 1/10^{-2} , 2/1 , 3/10^{2} , 4/10^{4} , 5/10^{6} , 6/10^{8}  }
	\def\listetrais{0,1,2,3,4,5,6}

	% --- Boucle sur toutes les positions ---
	\foreach \r in \listetrais {
    	% Initialisation
    	\global\def\found{0}
    	\xdef\nomtheta{}
    
   	 	% Vérification si \r est dans listetuple
    	\foreach \x/\y in \listetuple {
        	\ifdim \r pt=\x pt
            	\global\def\found{1}
            	\xdef\nomtheta{\y}
        	\fi
    	}

    	% Tracé de la graduation
    	\ifnum\found=1
        	% Graduation avec étiquette
        	\draw[color=\colorOne, thick, line width=0.5ex] 
            	(-0.1,\r) -- (0.1,\r) node[shift = {(-0,0)} , \colorOne , pos=-0.5 , left] {\large $\nomtheta$};
        	%\filldraw[line width=0.5ex, color=\colorSix] (\r,0) circle (4pt);
    		\else
        	% Graduation sans étiquette
        	\draw[color=\colorOne, thick, line width=0.5ex] 
            	(-0.1,\r) -- (0.1,\r);
        	%\filldraw[line width=0.5ex, color=\colorTwo] (\r,0) circle (4pt); 
    	\fi
	}

        
    %\drawgrid{-2}{10}{-1}{7};
\end{tikzpicture}


%\begin{tikzpicture}
%\begin{axis}[
%    xlabel={$\theta~(\mu\text{m}/\text{ms})$},
%    ylabel={$(\text{ms}/\mu\text{m}^2)$},
%    grid=both,
%    width=10cm,
%    height=6cm,
%    legend style={at={(0.98,0.98)}, anchor=north east, font=\small}
%]
%
%% rho_discr_c
%\addplot+[smooth, \colorSix, thick] 
%    table {rho_discr_c.dat};
%\addlegendentry{$\rho_{\rm ellipse}(\theta)$}
%
%% rho_vs_theta_T
%%\addplot+[smooth, \colorFour, thick] 
%%    table {../02_GGE_TBA/rho_vs_theta_T.dat};
%%\addlegendentry{$\rho(\theta)$}
%
%\end{axis}
%\end{tikzpicture}

\begin{tikzpicture}
\begin{axis}[
    xlabel={$\theta~(\mu\text{m}/\text{ms})$},
    ylabel={$\rho~(\text{ms}/\mu\text{m}^2)$},
    grid=both,
    width=10cm,
    height=6cm,
    legend style={at={(0.7,0.1)} , opacity = 0.8, anchor=south east, font=\small ,row sep=6pt  }
]



% rho T=100 nK
\addplot+[thick, \colorSix,  no marks , line width=0.5ex] table {../02_GGE_TBA/rho_vs_theta_T.dat};
\addlegendentry{\parbox{2cm}{$\gamma = 9\cdot 10^{-3}$ \\ $t = 86.5$}}


% rho T=0 nK
\addplot+[thick, \colorFour, no marks , line width=0.5ex] table {../02_GGE_TBA/rho_vs_theta_c.dat};
\addlegendentry{\parbox{2cm}{$\gamma = 9\cdot 10^{-3}$ \\ $t = 0$}}

% rho_s demi-cercle
\addplot+[ thick, \colorOne, dashed, no marks ,line width=0.5ex] table {../02_GGE_TBA/rho_discr_c.dat};
\addlegendentry{\parbox{2cm}{demi-cercle}}

\end{axis}
\end{tikzpicture}










\end{document}
